\chapter{Introduction}
\label{ch:introduction}

\section{Background}
\label{sec:background}

Electronic voting (e-voting) has emerged as a practical approach to modernizing democratic processes by replacing traditional paper-based balloting with digital systems. As nations worldwide seek to improve voter accessibility, reduce administrative costs, and accelerate result tabulation, e-voting systems have gained considerable attention from both the research community and government institutions. However, adopting electronic voting introduces core security challenges that do not exist in conventional voting: ballot privacy, vote integrity, coercion resistance, and universal verifiability must all be simultaneously guaranteed in a system where digital data can be copied, intercepted, or manipulated.

The integration of blockchain technology~\cite{nakamoto2008} with advanced cryptographic protocols offers a promising foundation for addressing these challenges. By leveraging the immutability and transparency properties of distributed ledgers alongside homomorphic encryption, zero-knowledge proofs, and anonymous credential schemes, it becomes possible to construct voting systems that provide mathematical guarantees of security rather than relying solely on organizational trust.

\subsection{Electronic Voting Systems}
\label{subsec:evoting-systems}

Electronic voting encompasses any voting method where the voter's intent is recorded through electronic means. These systems range from Direct Recording Electronic (DRE) machines used at physical polling stations to internet-based remote voting platforms that allow participation from any location.

\subsubsection{Evolution of electronic voting}
\label{subsubsec:evolution}

The evolution of electronic voting can be traced through several distinct generations. First-generation systems, deployed in the 1990s and early 2000s, were primarily DRE machines that digitized the ballot-marking process but stored votes in centralized databases controlled by election authorities. These systems improved tallying speed but introduced single points of failure and required voters to trust the machine vendor and election administrators completely.

Second-generation systems introduced cryptographic techniques such as mix-nets and blind signatures to provide some degree of voter anonymity. Notable examples include the Helios voting system, which pioneered web-based verifiable elections for organizational contexts. However, these systems relied on trusted third parties for key management and ballot processing, leaving them vulnerable to insider attacks and coercion.

Third-generation systems, emerging from 2018 onward, began integrating blockchain technology to provide immutable vote records and decentralized trust. Systems such as Voatz (used in US military overseas voting) and Agora (deployed as an observer in Sierra Leone's 2018 elections) represented early attempts at blockchain-based voting. However, these systems were criticized for weak cryptographic guarantees; a 2020 MIT security analysis of Voatz \cite{voatz_mit} revealed critical vulnerabilities including client-side attacks, server compromise possibilities, and vote manipulation risks.

Current e-voting research focuses on combining multiple cryptographic primitives into multi-layered security frameworks. The proposed framework represents this generation by integrating a multi-protocol cryptographic stack with a Caliper-benchmarked Hyperledger Fabric deployment to simultaneously address privacy, integrity, coercion resistance, and post-quantum security.

\subsubsection{E-voting system architecture}
\label{subsubsec:architecture}

A modern e-voting system architecture typically consists of four functional layers that interact to ensure secure vote processing:

\begin{enumerate}
    \item \textbf{Authentication Layer}: Responsible for verifying voter identity and eligibility. This may involve biometric verification, digital certificates, or credential-based authentication. The authentication mechanism must confirm the voter's right to participate without revealing their identity to other system components.
    
    \item \textbf{Computation Layer}: Handles the cryptographic operations that protect ballot privacy and ensure vote validity. This includes encryption, commitment generation, zero-knowledge proof construction, and digital signature creation. The computational demands of this layer directly impact the per-vote processing time and overall system throughput.
    
    \item \textbf{Privacy Layer}: Ensures that individual votes cannot be linked to specific voters even by system administrators. Privacy-preserving techniques include mix networks, homomorphic encryption, secret sharing, and anonymous aggregation protocols. The strength of this layer determines whether the system can resist surveillance and coercion attacks.
    
    \item \textbf{Storage Layer}: Provides persistent, tamper-evident storage of vote records. Blockchain technology serves this function by creating an append-only ledger that is distributed across multiple nodes, making retrospective vote manipulation computationally infeasible.
\end{enumerate}

\subsection{Security Challenges in E-Voting}
\label{subsec:security-challenges}

The security requirements of electronic voting are considerably more demanding than those of typical information systems. A voting system must simultaneously satisfy properties that are often in tension with one another; for example, votes must be both secret (unlinked to voters) and verifiable (provably counted correctly).

\subsubsection{Traditional security approaches}
\label{subsubsec:traditional-security}

Traditional approaches to e-voting security relied primarily on organizational controls and access restrictions. Votes were encrypted using standard symmetric or asymmetric algorithms (e.g., AES-256~\cite{nist_aes}, RSA~\cite{rivest1978rsa}), and system integrity depended on the trustworthiness of election authorities who held the decryption keys. This model suffers from several critical weaknesses: a compromised authority can decrypt and read individual votes, the tallying process is opaque to outside observers, and there is no mechanism to prove that votes were counted correctly without revealing them.

Beyond the tallying problem, traditional systems provide no protection against voter coercion. If a coercer can observe a voter's actions or demand proof of how they voted, the voter has no way to cast a free vote while appearing to comply with the coercer's demands. This vulnerability is particularly acute in remote voting scenarios where the voter is not protected by the secrecy of a polling booth.

\subsubsection{Cryptographic approaches}
\label{subsubsec:crypto-approaches}

Modern cryptographic approaches address these limitations through mathematical guarantees rather than organizational trust. Homomorphic encryption enables vote tallying without decrypting individual ballots. Zero-knowledge proofs allow voters to prove their votes are well-formed without revealing their content. Ring signatures provide anonymity within a group of potential signers while enabling double-vote detection through key images.

However, achieving all security properties simultaneously through cryptographic means requires careful protocol composition. Individual protocols provide specific guarantees, but their combination must be analyzed to ensure that the security properties are preserved under composition. This challenge of secure multi-protocol integration is a central concern of the proposed system design.

\subsection{Traditional E-Voting vs. Blockchain-Based E-Voting}
\label{subsec:traditional-vs-blockchain}

The key distinction between traditional and blockchain-based e-voting lies in the trust model. Traditional systems require voters to trust a centralized authority, while blockchain-based systems distribute trust across a network of independent nodes.

\subsubsection{Core differences}
\label{subsubsec:core-diff}

In traditional e-voting, a central election authority controls the encryption keys, manages the vote database, and performs the tally. Voters must trust that this authority will not manipulate votes, leak individual ballots, or selectively exclude certain votes from the count. The verification process, if any, is typically limited to auditing by authorized observers who must again be trusted.

Blockchain-based e-voting fundamentally alters this trust model. Votes are recorded on a distributed ledger that is maintained by multiple independent parties. The immutability property of the blockchain ensures that once a vote is recorded, it cannot be modified or deleted without detection. Smart contracts can enforce voting rules automatically, removing the need to trust election officials to follow procedures correctly. Cryptographic proofs attached to each vote enable universal verifiability: any observer can verify that all recorded votes are well-formed and correctly tallied.

\subsubsection{Advantages of blockchain-based approach}
\label{subsubsec:blockchain-advantages}

Several advantages of blockchain-based e-voting are worth highlighting. First, the decentralized nature of the blockchain eliminates single points of failure: no single entity can compromise the entire election. Second, the immutable ledger provides an auditable trail that persists indefinitely, enabling post-election verification at any time. Third, smart contracts ensure consistent enforcement of election rules without human intervention. Fourth, blockchain transparency enables end-to-end verifiability, where voters can verify that their votes were recorded and counted correctly.

These properties are particularly relevant in contexts where institutional trust is low, such as elections in developing nations, organizational governance, or scenarios where election manipulation is a concern.

\subsection{Blockchain Technology in E-Voting}
\label{subsec:blockchain-tech}

Blockchain platforms used for e-voting fall into two broad categories: public blockchains (e.g., Ethereum~\cite{wood2014ethereum}) and permissioned blockchains (e.g., Hyperledger Fabric~\cite{androulaki2018}). Public blockchains offer maximum decentralization but suffer from scalability limitations, high transaction costs, and energy consumption concerns. Permissioned blockchains provide controlled access, higher throughput, and configurable privacy, making them more suitable for election systems where the set of validating nodes is known and controlled by election authorities.

Our framework employs Hyperledger Fabric v2.5~\cite{androulaki2018}, a permissioned blockchain framework maintained by the Linux Foundation. Fabric's modular architecture supports pluggable consensus mechanisms, channel-based data isolation, and chaincode (smart contracts) written in general-purpose languages such as Go. The Fabric Gateway SDK enables application-level interaction with the blockchain over gRPC channels secured with TLS, providing authenticated communication between the voting application and the distributed ledger. Recent empirical studies~\cite{fabric_perf} have demonstrated that Fabric achieves substantially higher throughput and lower latency than public blockchain alternatives for enterprise applications.

\subsection{Cryptographic Protocols for E-Voting}
\label{subsec:crypto-protocols}

The security of the proposed framework rests on a multi-protocol cryptographic stack in which each primitive provides a specific security guarantee:

\begin{itemize}
    \item \textbf{Paillier Homomorphic Encryption}~\cite{paillier1999}: Enables encrypted vote tallying without decryption. Based on the Decisional Composite Residuosity Assumption (DCRA)~\cite{paillier1999}, Paillier's additive homomorphism allows the election result to be computed directly from encrypted ballots: $E(a) \times E(b) = E(a + b)$.
    
    \item \textbf{Pedersen Commitments}: Provide perfectly hiding and computationally binding commitments to vote values. Used as the foundation for zero-knowledge proofs that verify vote validity.
    
    \item \textbf{Zero-Knowledge Proofs ($\Sigma$-Protocol)}: Enable voters to prove that their encrypted votes contain valid values ($w \in \{0, 1\}$) and sum to exactly one ($\sum w = 1$) without revealing which candidate received the vote. The Strong Fiat-Shamir transformation~\cite{bernhard2012} is used to make the proofs non-interactive, with the public parameters, statement, election identifier, and a per-proof nonce bound into the challenge to prevent the malleability attack identified by Bernhard, Pereira, and Warinschi against earlier non-interactive Fiat-Shamir constructions used in voting systems. In the deployed pipeline, both the per-slot binary proofs and the sum-one proof are verified at credential issuance and re-checked at the time of vote casting, so the soundness reduction stated in Chapter~\ref{ch:methodology} as Theorem~2 (Vote Validity) is enforced at every entry point that consumes a credential.
    
    \item \textbf{Linkable Ring Signatures}: Provide voter anonymity within a ring of 100 members while enabling double-vote detection through deterministic key images. Based on the Liu-Wei-Wong construction \cite{liu2004}.
    
    \item \textbf{SMDC (Self-Masking Deniable Credentials)}: A novel coercion-resistance mechanism in which each voter receives $k=5$ credential slots (1 real + 4 fake). The slot commitments are perfectly hiding under Pedersen, and the real-slot index itself is derived from a server-held HMAC secret $k_{\text{srv}}$ that an external observer cannot recompute from public inputs. Under coercion, voters can present any fake slot to produce valid-looking but silently discarded votes. The index is disclosed to the legitimate voter only through a duress-gated reveal endpoint.
    
    \item \textbf{SA\textsuperscript{2} (Samplable Anonymous Aggregation)}: A two-server privacy model adapted from Apple's federated learning research \cite{sa2_apple}. Votes are split into masked shares distributed between two non-colluding servers. As long as one server remains honest, individual votes cannot be reconstructed.
    
    \item \textbf{Kyber-768 Post-Quantum Encryption}: NIST Level 3 post-quantum key encapsulation based on Module-LWE \cite{nist_kyber}, providing protection against future quantum computer attacks through hybrid encryption with Paillier.

    \item \textbf{DJN Threshold Paillier Decryption}~\cite{damgard2010}: Distributes the Paillier private key across $n$ trustees via Shamir's $(t,n)$-secret sharing~\cite{shamir1979}, so that no single entity can unilaterally decrypt election results; any $t$ trustees jointly reconstruct the tally through Lagrange interpolation.
\end{itemize}

\subsection{State of the Art in the Modern Context}
\label{subsec:state-of-art}

Recent blockchain e-voting research has produced systems that each strengthen one or two security properties but none that simultaneously achieve all of: homomorphic tallying for exact results, formal zero-knowledge proofs for vote validity, linkable ring signatures for anonymity with double-vote detection, deniable credentials for coercion resistance, two-server anonymous aggregation, post-quantum security, and a Caliper-benchmarked Hyperledger Fabric deployment with empirical performance data. Recurring shortcomings across the surveyed work include super-linear tally complexity, reliance on trusted counting authorities or differential-privacy approximations, weak symmetric-only privacy, and the absence of a coercion-resistance mechanism that survives a coercer who can demand proof of how the voter cast a ballot. A detailed analysis of representative works, their reported performance, and the precise gaps they leave open is presented in Chapter~\ref{ch:related-works}. The proposed framework directly addresses these gaps and is positioned in detail against the cited systems in Section~\ref{sec:comparative} of Chapter~\ref{ch:results}.

\subsection{Importance and Relevance of the Topic}
\label{subsec:importance}

Secure electronic voting matters beyond technical innovation; it touches on democratic principles directly. In many developing nations, voter intimidation and coercion remain serious obstacles to free elections. International election observers frequently report incidents of vote buying, family voting, and employer pressure on voting choices. A voting system that provides mathematical guarantees of coercion resistance, not just procedural safeguards, would mark a genuine step forward for democratic integrity.

On a related note, the emergence of quantum computing poses a long-term threat to current cryptographic systems~\cite{shor1997}. Elections may produce records that must remain confidential for decades. A vote encrypted today with RSA, ECC, or any scheme whose security reduces to integer factorisation or discrete logarithms could potentially be decrypted by a future quantum computer running Shor's algorithm~\cite{shor1997}, retroactively compromising voter privacy. This \emph{harvest-now-decrypt-later} threat model has motivated the U.S. National Institute of Standards and Technology to standardise lattice-based replacements, including ML-KEM (FIPS~203~\cite{nist_kyber}) for key encapsulation and ML-DSA (FIPS~204~\cite{nist_mldsa}) for digital signatures, both finalised in 2024. Post-quantum cryptographic protection is therefore not merely a theoretical concern but a practical and standards-aligned necessity for long-term ballot secrecy.

From a scalability perspective, nations with large populations require voting systems whose per-candidate tally cost grows only linearly with the number of voters $N$ and whose total tally cost grows no worse than linearly with the number of candidates $m$. The proposed framework achieves $O(N)$ per candidate and $O(N \cdot m)$ total tally complexity, confirmed through empirical benchmarks --- a factor-$m$ improvement over ISE-Voting's $O(N \cdot m^{2})$. Feasibility for national-scale deployments is further evidenced by approximately 200~TPS Caliper-verified chaincode throughput and 100\% tally integrity across thirteen end-to-end election simulations spanning 100 to 100{,}000 voters.

The intersection of blockchain technology, advanced cryptography, and post-quantum security in the context of electronic voting represents an active and important area of research with direct societal impact.

% ============================================================
\section{Motivation}
\label{sec:motivation}

The development of the proposed framework is motivated by several gaps observed in the blockchain-based e-voting literature surveyed in Chapter~\ref{ch:related-works}. Substantial progress has been made on each individual security property, but the surveyed work leaves room for a more tightly integrated composition.

First, the works surveyed in Chapter~\ref{ch:related-works} typically address only a subset of the required security properties. Schemes built around homomorphic tallying often omit a deployable coercion-resistance mechanism, while the seminal coercion-resistance constructions of JCJ~\cite{jcj2005} and its instantiation in Civitas~\cite{civitas2008} introduced the deniable-credential paradigm and were realised as research prototypes. Building on that line of work, the present thesis integrates a JCJ-style deniable credential mechanism (SMDC) into a Hyperledger Fabric permissioned-blockchain pipeline that also incorporates post-quantum hybrid encryption, two-server anonymous aggregation, and linkable ring signatures --- a composition that, to the best of our knowledge and within the works surveyed in Chapter~\ref{ch:related-works}, has not been previously reported.

Second, several of the surveyed tallying mechanisms exhibit super-linear scaling in the number of candidates. ISE-Voting's $O(N \times m^2)$ complexity, for example, requires approximately $10^8$ operations for a one-million-voter election with ten candidates, motivating a linear-in-voters and linear-in-candidates tally for national-scale deployments.

Third, many of the systems reviewed in Chapter~\ref{ch:related-works} report theoretical designs or local-only simulations; an end-to-end Hyperledger Fabric integration with empirical throughput and latency measurements on a permissioned platform is reported in only a subset of these works. Bridging this design--deployment gap with reproducible Hyperledger Caliper benchmarks is one of the engineering objectives of this thesis.

Fourth, although standalone post-quantum voting proposals exist in the broader literature, the integration of post-quantum key encapsulation into the same pipeline as coercion-resistant credentials, ring-signature anonymity, and two-server anonymous aggregation is not reported among the works surveyed in Chapter~\ref{ch:related-works}. The proposed framework adds a Kyber-768 hybrid encryption layer to address the harvest-now-decrypt-later threat described in Section~\ref{subsec:importance}.

These observations motivate the design and implementation of the proposed framework as an integrated e-voting system that composes a multi-protocol cryptographic stack on real permissioned-blockchain infrastructure, with the cast endpoint wired end-to-end through the Fabric Gateway SDK so that successfully cast votes are published to the deployed chaincode during normal operation.

% ============================================================
\section{Scope and Limitations}
\label{sec:scope}

This thesis focuses on the cryptographic protocol design, implementation, and performance evaluation of the proposed system. The scope and boundaries of this research are defined as follows:

\subsection{In Scope}
\begin{itemize}
    \item Design and implementation of a multi-protocol cryptographic stack (Paillier homomorphic encryption~\cite{paillier1999}, Pedersen commitments~\cite{pedersen1991}, zero-knowledge $\Sigma$-proofs with Strong Fiat-Shamir~\cite{bernhard2012}, linkable ring signatures~\cite{liu2004}, the proposed SMDC deniable credentials, two-server SA\textsuperscript{2} aggregation~\cite{sa2_apple}, Kyber-768 (CRYSTALS-Kyber round-3 via CIRCL~\cite{nist_kyber}), and Damg\aa rd-Jurik-Nielsen threshold Paillier decryption~\cite{damgard2010}) in Go~1.24.
    \item Integration with Hyperledger Fabric~v2.5~\cite{androulaki2018} as the underlying permissioned-blockchain infrastructure.
    \item Empirical performance evaluation including crypto micro-benchmarks, pipeline analysis, scalability testing (100 to 100,000 voters), and Caliper blockchain benchmarks (10,000 to 100,000 transactions).
    \item Security analysis with formal theorem statements, proof sketches for six security properties, and ProVerif formal verification of three core properties.
    \item Comparative analysis with five state-of-the-art blockchain e-voting systems.
\end{itemize}

\subsection{Out of Scope}
The following items fall outside the boundary of the present work: user interface (UI/UX) design and implementation for the voter-facing application, voter registration infrastructure and national identity integration, formal verification using Tamarin or other automated tools beyond the ProVerif symbolic model, multi-constituency deployment with geographic distribution, mobile client implementation and accessibility testing, and real-world election deployment and usability studies with actual voters. These out-of-scope items are discussed as future work directions in Chapter~\ref{ch:conclusion}.

% ============================================================
\section{Problem Statement}
\label{sec:problem-statement}

Despite the advances in blockchain-based e-voting, significant challenges remain in building a system that provides comprehensive security guarantees while maintaining practical performance. The core problems addressed by this thesis are as follows:

\begin{itemize}
    \item How can multiple cryptographic protocols (homomorphic encryption, zero-knowledge proofs, ring signatures, deniable credentials, anonymous aggregation, and post-quantum key encapsulation) be securely composed into a unified voting pipeline that preserves the individual security properties of each protocol under composition?
    
    \item How can coercion resistance be practically implemented in a blockchain e-voting context, such that a coerced voter can present plausible fake credentials to a coercer without the coercer being able to distinguish real from fake votes, while ensuring that only real votes are counted in the final tally?
    
    \item How can the vote tallying process achieve $O(N)$ computational complexity \emph{per candidate} (linear in the voter count $N$), with $O(N \cdot m)$ total cost across $m$ candidates, while preserving the exact correctness of results without introducing noise or approximation?

    \item How can a deployed Hyperledger Fabric blockchain be integrated with the cryptographic voting pipeline to provide immutable vote records and chaincode-level verification, and what throughput and latency characteristics does this integration achieve under realistic load conditions?
\end{itemize}

% ============================================================
\section{Research Objectives}
\label{sec:research-objectives}

The primary goal of this research is to design, implement, and empirically evaluate a blockchain-based e-voting system that provides comprehensive cryptographic security guarantees while demonstrating practical performance on a deployed Hyperledger Fabric network. The specific objectives are:

\begin{itemize}
    \item To design and implement a multi-protocol cryptographic voting pipeline that provides ballot privacy (Paillier homomorphic encryption), vote validity (Pedersen commitments with zero-knowledge $\Sigma$-proofs), voter anonymity (linkable ring signatures), coercion resistance (SMDC deniable credentials), privacy-preserving aggregation (SA\textsuperscript{2} two-server model), post-quantum security (Kyber-768 hybrid encryption), and distributed key management (Damg\aa rd-Jurik-Nielsen threshold Paillier decryption).

    \item To develop a novel coercion resistance mechanism (SMDC) based on deniable credentials with $k=5$ slots, integrated with behavioural duress detection, that provides computationally indistinguishable real and fake credentials in a deployed blockchain voting context.

    \item To achieve $O(N)$ per-candidate tally complexity (and $O(N \cdot m)$ total) through Paillier's additive homomorphism, enabling linear-in-voters vote counting that grows only linearly, not quadratically, with the number of candidates, and to validate this complexity through empirical benchmarks from 100 to 100{,}000 voters.

    \item To integrate the voting system with a deployed Hyperledger Fabric v2.5 network and empirically measure throughput and latency using Hyperledger Caliper benchmarks across transaction volumes from 10{,}000 to 100{,}000.
\end{itemize}

% ============================================================
\section{Research Contributions}
\label{sec:contributions}

This thesis makes the following contributions to the field of blockchain-based electronic voting. To avoid overstating novelty, the contributions are explicitly grouped into (a) genuinely novel cryptographic constructions, (b) system-integration and engineering contributions, and (c) empirical results.

\begin{enumerate}
    \item \textbf{SMDC Coercion Resistance with Duress-Bound Index Disclosure (Novel Construction).} The thesis introduces Self-Masking Deniable Credentials (SMDC), a practical coercion-resistance mechanism in which each voter receives $k=5$ credential slots (1 real and 4 fake). The real-slot index is derived via $\text{HMAC-SHA256}\bigl(k_{\text{srv}},\ \text{``smdc-real-index:''} \,\Vert\, \text{vid} \,\Vert\, \text{eid}\bigr) \bmod k$ where $k_{\text{srv}}$ is a server-held secret kept confidential to the election authority. The index is therefore re-derivable only by the authority and is \emph{never} returned in the credential-issuance response. To learn the real index, a voter must submit the duress signal registered at issue time to a dedicated reveal endpoint that returns the index only on a constant-time signal match; a registration-time or post-registration coercer who knows only the public pair $(\text{vid}, \text{eid})$ therefore obtains nothing exploitable, and a voter under duress can deceive the coercer by presenting any wrong signal. SMDC delivers a complete implementation of deniable credentials integrated with a blockchain e-voting pipeline, addressing the implementation gap left open by the theoretical credential-management schemes of JCJ~\cite{jcj2005} and Civitas~\cite{civitas2008} for the works surveyed in Chapter~\ref{ch:related-works}.

    \item \textbf{Behavioural Duress Detection (Novel Construction).} A behavioural coercion-detection mechanism is presented, using HMAC-SHA256 hashing of a voter-chosen behavioural signal (blink count, head tilt, long press, time delay, voice command, etc.) under a server-held key $K_d$, with constant-time equality verification against a per-voter stored hash $h_{\text{stored}}$ that is recorded at registration time. The same signal gates the SMDC reveal-slot-index endpoint (above) and the cast-time coercion short-circuit (below), so a single registered secret simultaneously protects (i)~the voter's ability to retrieve the real index without leaking it to a coercer at registration and (ii)~the voter's ability to silently discard a coerced vote at cast time. The deterministic per-voter equality check, the registration-time signal capture, and the dual-purpose gating are original to this work.

    \item \textbf{Coercion Short-Circuit Semantics (Novel Construction).} The vote-casting pipeline is extended with silent-substitution semantics gated on the behavioural-duress equality check: when the presented signal fails to match $h_{\text{stored}}$, the behavioural weight is zeroed but the full ring-signature and SA\textsuperscript{2} share-split steps still execute (preserving timing indistinguishability), and the per-candidate ciphertexts are scalar-multiplied by zero so that the homomorphic tally counts $E(0)$. Crucially, the key image is \emph{not} marked used in this path, so the voter retains the ability to cast a genuine vote once the coercer is gone. The receipt returned to a coerced voter is observationally identical to a genuine-cast receipt. This short-circuit triggers on the behavioural signal only and is orthogonal to the SMDC slot choice (a genuine vote on a fake SMDC slot still proceeds to chain as a legitimate decoy).

    \item \textbf{Multi-Protocol Pipeline Integration (Engineering Contribution).} This work presents a unified pipeline that composes the multi-protocol cryptographic stack --- Paillier homomorphic encryption~\cite{paillier1999}, Pedersen commitments~\cite{pedersen1991}, zero-knowledge $\Sigma$-proofs (Strong Fiat-Shamir)~\cite{bernhard2012}, linkable ring signatures~\cite{liu2004}, SA\textsuperscript{2} anonymous aggregation~\cite{sa2_apple}, Kyber-768 post-quantum hybrid encryption~\cite{nist_kyber}, and DJN threshold Paillier decryption~\cite{damgard2010} --- together with the novel SMDC and duress mechanisms above into a single voting pipeline. None of the works surveyed in Chapter~\ref{ch:related-works} reports the same combination of primitives in a single integrated pipeline. The underlying primitives are existing constructions; the contribution lies in their composition, the resolution of inter-protocol parameter compatibility, and the deployment engineering.

    \item \textbf{SA\textsuperscript{2} Adaptation for Voting (Integration Choice, not Novel).} Apple's Samplable Anonymous Aggregation primitive~\cite{sa2_apple}, originally designed for federated data analysis, is integrated with Paillier additive homomorphism so that the encrypted tally is reconstructed via mask cancellation across two non-colluding aggregators. The SA\textsuperscript{2} construction itself is not a contribution of this thesis; the contribution is its placement in the voting pipeline.

    \item \textbf{Hyperledger Fabric Integration (Engineering Contribution).} The vote-casting pipeline is wired end-to-end through the Fabric Gateway SDK (fabric-gateway v1.10.1): every non-coerced vote produced by the API server's cast endpoint is published to the deployed chaincode via \texttt{SubmitTransaction}, and the resulting chain transaction identifier is returned in the voter's receipt. Independent Hyperledger Caliper benchmarks at the chaincode level report approximately 200~TPS throughput with a 100\% success rate at 10\,000 transactions; an end-to-end Caliper run that exercises the integrated cast endpoint with a live Fabric peer in the loop is identified as future work.

    \item \textbf{$O(N)$-per-Candidate Tally with Empirical Validation (Result).} Through Paillier's additive homomorphism, the tally achieves $O(N)$ complexity per candidate and $O(N \cdot m)$ total complexity (with $N$ denoting the voter count, distinct from the Paillier modulus $n$), confirmed empirically: the per-vote per-candidate cost is approximately 6.1~$\mu$s across four orders of magnitude (100 to 100{,}000 voters), and the total tally time scales linearly with $m$ (Section~\ref{subsubsec:candidate-scale}). This is a factor-$m$ improvement over ISE-Voting's $O(N \cdot m^{2})$, and the 100\% exact tally contrasts with BP-Vot's $\geq 98\%$ differential-privacy approximation.
\end{enumerate}

% ============================================================
\section{Thesis Outline}
\label{sec:outline}

The remainder of this thesis is organized as follows:

\begin{itemize}
    \item \textbf{Chapter 2, Related Works}, discusses the existing research on blockchain-based e-voting systems with a focus on security properties, cryptographic techniques, and coercion resistance mechanisms. It reviews six key works in detail, analyzing their methodologies, contributions, and limitations. The chapter concludes with a comprehensive comparison table that positions our framework relative to the state of the art.
    
    \item \textbf{Chapter 3, System Design and Cryptographic Protocols}, presents the complete architecture of the proposed system. This chapter includes the detailed mathematical formulation of the multi-protocol cryptographic stack, the twenty-step vote casting pipeline, the SMDC credential generation mechanism, the SA\textsuperscript{2} aggregation protocol, and the blockchain integration architecture. Engineering and societal considerations are also addressed.
    
    \item \textbf{Chapter 4, Implementation and Performance Evaluation}, describes the implementation details including the technology stack (Go 1.24, Gin v1.11, Hyperledger Fabric v2.5.15), test coverage analysis, and comprehensive performance benchmarks. Results include crypto micro-benchmarks, end-to-end pipeline measurements, scalability validation from 100 to 100{,}000 voters, and Hyperledger Caliper blockchain throughput benchmarks from 10{,}000 to 100{,}000 transactions. Comparative analysis with existing systems is presented.
    
    \item \textbf{Chapter 5, Conclusion and Future Work}, summarizes the key findings of this research, discusses the system's limitations, and outlines directions for future work including formal verification, lattice-based ring signatures, and production deployment considerations.
\end{itemize}

% ============================================================
\section{Conclusion}
\label{sec:intro-conclusion}

This chapter has provided the background, motivation, and problem statement for this research. The fundamental security challenges in electronic voting (ballot privacy, vote integrity, coercion resistance, and post-quantum security) were identified, along with the limitations of existing blockchain-based solutions. The research objectives were defined to address these gaps through a multi-protocol cryptographic stack deployed on a Caliper-benchmarked Hyperledger Fabric network. The next chapter reviews the existing literature in detail, analyzing the methodology, contributions, and limitations of state-of-the-art blockchain e-voting systems.
